\documentclass[../../main.tex]{subfiles}
\begin{document}
\section{习题}\label{section:ch21-problems}

\begin{exercise}\label{exercise:ch21-1}
设 Lie 群 $G$ 连续作用于流形 $M$。证明：若定义该作用的映射 $G\times M\to M$ 是真映射，则该作用是真作用。举反例说明反之不必成立。
\end{exercise}

\begin{exercise}\label{exercise:ch21-2}
令 $X=\mathbb{R}^n/\mathbb{R}^+$，其中 $\mathbb{R}^+$ 的作用由 (21.1) 给出。证明 $X$ 有开子集同胚于 $\mathbb{S}^{n-1}$，并有一点属于每个非空闭子集。
\end{exercise}

\begin{exercise}\label{exercise:ch21-3}
证明 $\mathbb{S}^1$ 在 $\mathbb{S}^{2n+1}$ 上的 Hopf 作用（习题 7-11）自由且真，且轨道空间 $\mathbb{S}^{2n+1}/\mathbb{S}^1$ 微分同胚于 $\mathbb{CP}^n$。（提示：考虑自然商映射 $\mathbb{C}^{n+1}\setminus\{0\}\to\mathbb{CP}^n$ 在 $\mathbb{S}^{2n+1}$ 上的限制。商映射 $\mathbb{S}^{2n+1}\to\mathbb{CP}^n$ 称为 Hopf 映射。）
\end{exercise}

\begin{exercise}\label{exercise:ch21-4}
设 $M$ 是光滑 $n$-流形，$V$ 是 $M$ 上光滑向量场，使 $V$ 的每条积分曲线都以同一周期为周期（习题 9-1）。在 $M$ 上定义等价关系 $p\sim q$：若 $p,q$ 在 $V$ 同一条积分曲线的像中。令 $M/\sim$ 为商空间，$\pi\colon M\to M/\sim$ 为商映射。证明 $M/\sim$ 是拓扑 $(n-1)$-流形，并有唯一光滑结构使 $\pi$ 为光滑淹没。
\end{exercise}

\begin{exercise}\label{exercise:ch21-5}
证明商流形定理的下列部分逆：若 Lie 群 $G$ 光滑自由作用于光滑流形 $M$，且轨道空间 $M/G$ 有光滑流形结构使商映射 $\pi\colon M\to M/G$ 为光滑淹没，则 $G$ 的作用真。（提示：先证对 $\pi$ 的任意光滑局部截面 $\sigma\colon U\to M$，映射 $(g,x)\mapsto g\cdot\sigma(x)$ 是从 $G\times U$ 到 $\pi^{-1}(U)$ 的微分同胚。）
\end{exercise}

\begin{exercise}\label{exercise:ch21-6}
设 Lie 群 $G$ 光滑、自由、真作用于光滑流形 $M$。证明 $M$ 是以 $M/G$ 为底、以 $G$ 为模型纤维、以商映射 $\pi\colon M\to M/G$ 为投影的光滑纤维丛的全空间。（提示：见习题 21-5 的提示。）由习题 10-19 推出 $M$ 紧当且仅当 $G$ 与 $M/G$ 都紧。（注：这样得到的纤维丛称为\emph{主 $G$-丛}。）
\end{exercise}

\begin{exercise}\label{exercise:ch21-7}
设离散 Lie 群 $\Gamma$ 光滑、自由、真作用于连通光滑流形 $M$。证明 $M$ 上的 Riemann 度量 $g$ 是度量在 $M/\Gamma$ 上经商映射 $\pi\colon M\to M/\Gamma$ 的拉回，当且仅当 $\Gamma$ 以 $g$ 的等距作用（即对每个 $\gamma\in\Gamma$，$x\mapsto\gamma\cdot x$ 是等距）。
\end{exercise}

\begin{exercise}\label{exercise:ch21-8}
设 $M$ 是光滑流形，$\pi\colon E\to M$ 是 $M$ 上光滑向量丛。设离散 Lie 群 $\Gamma$ 在 $E$ 与 $M$ 上都光滑、自由、真作用，且 $\pi$ 是 $\Gamma$-等变的，对每个 $p\in M$、$g\in\Gamma$，由 $v\mapsto g\cdot v$ 给出的映射 $E_p\to E_{g\cdot p}$ 线性。证明 $E/\Gamma$ 可赋以 $M/\Gamma$ 上光滑向量丛的结构，使图表交换：$E\to E/\Gamma$，$M\to M/\Gamma$。（水平映射为商映射，竖直为丛投影。）
\end{exercise}

\begin{exercise}\label{exercise:ch21-9}
证明 $\mathbb{Z}$ 在 $\mathbb{R}^2$ 上由 $n\cdot(x,y)=(x+n,(-1)^n y)$ 给出的作用光滑、自由、真。结合习题 21-8 的结果，给出例 10.3 的 Möbius 丛是 $\mathbb{S}^1$ 上非平凡光滑秩 1 向量丛的简单证明，并证明用于定义 $E$ 的商映射 $q\colon\mathbb{R}^2\to E$ 是光滑正规覆叠映射，其覆叠自同构群即上述 $\mathbb{Z}$ 在 $\mathbb{R}^2$ 上的作用。
\end{exercise}

\begin{exercise}\label{exercise:ch21-10}
令 $\mathbb{CP}^n$ 表示 $n$ 维复射影空间（习题 1-9）。证明 $U(n+1)$ 在 $\mathbb{C}^{n+1}\setminus\{0\}$ 上的自然作用下降为 $\mathbb{CP}^n$ 上的光滑可迁作用，故 $\mathbb{CP}^n$ 是齐性 $U(n+1)$-空间。取一点并确定其迷向群。
\end{exercise}

\begin{exercise}\label{exercise:ch21-11}
本题证明非连通光滑流形是齐性的，当且仅当其各分支彼此微分同胚且其中一个齐性。
(a) 设 $M$ 是齐性 $G$-空间。证明其各分支彼此微分同胚，且存在开子群 $G_0\subseteq G$，使每个分支是齐性 $G_0$-空间。
(b) 反之，设 $M$ 是光滑流形，其各分支彼此微分同胚，且有 Lie 群 $G_0$ 光滑可迁作用于其中一个分支 $M_0$。证明直积群 $G_0\times\mathbb{Z}$ 光滑可迁作用于 $M$。（提示：先证 $M$ 微分同胚于 $M_0$ 与可数离散空间的乘积。）
\end{exercise}

\begin{exercise}\label{exercise:ch21-12}
证明 Grassmann 流形 $G_k(\mathbb{R}^n)$ 上由例 21.21 定义的光滑结构与例 1.36 定义的相同。
\end{exercise}

\begin{exercise}\label{exercise:ch21-13}
设 $V$ 是有限维实向量空间。证明对每个 $k$，Grassmann 流形 $G_k(V)$ 紧。（提示：证明它是某紧 Lie 群的商空间。）
\end{exercise}

\begin{exercise}\label{exercise:ch21-14}
设 $V$ 是 $n$ 维实向量空间，$0<k<n$，$G_k(V)$ 为 $k$ 维子空间的 Grassmann 流形。令 $\mathbb{P}(\wedge^k V)$ 表示 $\wedge^k V$ 的射影化（习题 2-11）。定义 $\iota\colon G_k(V)\to\mathbb{P}(\wedge^k V)$，$\iota(S)=[v_1\wedge\cdots\wedge v_k]$（$S=\operatorname{span}(v_1,\dots,v_k)$）。证明 $\iota$ 良定义，且是光滑嵌入，其像为 $\wedge^k V$ 中所有非零可分解元的等价类集合。（称为 Plücker 嵌入。）
\end{exercise}

\begin{exercise}\label{exercise:ch21-15}
$\mathbb{C}^n$ 的 $k$ 维复线性子空间集合记为 $G_k(\mathbb{C}^n)$。证明 $G_k(\mathbb{C}^n)$ 有唯一光滑流形结构使其成为齐性 $\operatorname{GL}(n,\mathbb{C})$-空间，作用由 $\operatorname{GL}(n,\mathbb{C})$ 在 $\mathbb{C}^n$ 上通常作用诱导。证明 $G_k(\mathbb{C}^n)$ 紧并确定其维数。
\end{exercise}

\begin{exercise}\label{exercise:ch21-16}
设 $F_K(V)$ 为例 21.22 中有限维向量空间 $V$ 中类型 $K$ 的旗集合。证明 $\operatorname{GL}(V)$ 在 $F_K(V)$ 上可迁，某特定旗的迷向群是 $\operatorname{GL}(V)$ 的闭子群，从而 $F_K(V)$ 有唯一光滑流形结构使该作用光滑。$\dim F_K(V)$ 是多少？对哪些 $K$，$F_K(V)$ 紧？
\end{exercise}

\begin{exercise}\label{exercise:ch21-17}
设 Lie 群光滑（不必真或自由）作用于光滑流形 $M$。证明每条轨道是 $M$ 的浸入子流形；若作用真，则它是嵌入的。
\end{exercise}

\begin{exercise}\label{exercise:ch21-18}
群 $G$ 的\emph{中心}是与 $G$ 中所有元素交换的元素集合；$G$ 的子群称为\emph{中心的}，若它含于 $G$ 的中心。证明连通 Lie 群的每个离散正规子群是中心的。（提示：用习题 7-8 的结果。）
\end{exercise}

\begin{exercise}\label{exercise:ch21-19}
用定理 7.7 与习题 21-18 的结果证明每个 Lie 群的基本群是交换群。可免证明引用：若 $\pi\colon\widetilde G\to G$ 是万有覆叠，则自同构群 $\operatorname{Aut}_\pi(\widetilde G)$ 同构于 $\pi_1(G,e)$。
\end{exercise}

\begin{exercise}\label{exercise:ch21-20}
(a) 设 $G$ 是连通交换 Lie 群。证明 $G$ 的指数映射是从 $\operatorname{Lie}(G)$（按向量加法视为 Lie 群）到 $G$ 的满射 Lie 群同态。
(b) 证明每个连通交换 Lie 群同构于 $\mathbb{R}^k\times\mathbb{T}^l$（某非负整数 $k,l$）。
\end{exercise}

\begin{exercise}\label{exercise:ch21-21}
证明 $\operatorname{SO}(3)$ 同构于 $\operatorname{SU}(2)/\{\pm I_2\}$，且微分同胚于 $\mathbb{RP}^3$，如下：
(a) 令 $\mathbb{S}^3\subseteq\mathbb{H}$ 为单位四元数群，$E\subseteq\mathbb{H}$ 为虚四元数子空间（习题 7-22、7-23、8-6）。对 $q\in\mathbb{S}^3$，证明由 $v\mapsto qvq^{-1}$ 给出的线性映射 $\mathbb{H}\to\mathbb{H}$ 把 $E$ 映入自身，并保持 $E$ 上内积 $\langle v,w\rangle=\frac12(v^\ast w+w^\ast v)$。
(b) 对 $q\in\mathbb{S}^3$，令 $\rho(q)$ 为映射 $v\mapsto qvq^{-1}$ 关于 $E$ 的基 $(i,j,k)$ 的矩阵表示。证明 $\rho(q)\in\operatorname{SO}(3)$，且 $\rho\colon\mathbb{S}^3\to\operatorname{SO}(3)$ 是满射 Lie 群同态，核为 $\{\pm1\}$。
(c) 证得结论。
\end{exercise}

\begin{exercise}\label{exercise:ch21-22}
若 $G,H$ 是 Lie 群，且存在从 $G$ 到 $H$、核为 $G_0$ 的满射 Lie 群同态，则称 $G$ 是 $G_0$ 经 $H$ 的\emph{扩张}。令 $\mathfrak g$ 为任意有限维 Lie 代数，证明 Lie 代数同构于 $\mathfrak g$ 的非连通 Lie 群恰为连通群经离散 Lie 群的扩张。（提示：见命题 7.15。用于 p.557。）
\end{exercise}

\begin{exercise}\label{exercise:ch21-23}
证明 $\operatorname{GL}^+(n,\mathbb{R})$ 微分同胚于 $\operatorname{SO}(n)\times\mathbb{R}^{n(n+1)/2}$。（提示：用命题 21.35 的 QR 分解构造微分同胚 $\operatorname{SO}(n)\times\operatorname{TC}(n,\mathbb{R})\cong\operatorname{GL}^+(n,\mathbb{R})$，其中 $\operatorname{TC}(n,\mathbb{R})$ 为对角元正的 $n\times n$ 上三角实矩阵组成的 Lie 群。）
\end{exercise}

\begin{exercise}\label{exercise:ch21-24}
证明 $\operatorname{GL}(n,\mathbb{C})$ 微分同胚于 $U(n)\times\mathbb{R}^{n^2}$，从而连通。（提示：仿习题 21-23，但用 Hermite 点积 $z\cdot w=\sum_i z_i\overline{w_i}$ 代替 Euclid 点积。）
\end{exercise}

\begin{exercise}\label{exercise:ch21-25}
证明 $\operatorname{SL}(n,\mathbb{R})$、$\operatorname{SL}(n,\mathbb{C})$ 分别微分同胚于 $\operatorname{SO}(n)\times\mathbb{R}^{n(n+1)/2-1}$、$\operatorname{SU}(n)\times\mathbb{R}^{n^2-1}$。推出 $\operatorname{SL}(n,\mathbb{R})$、$\operatorname{SL}(n,\mathbb{C})$ 连通，$\operatorname{SL}(2,\mathbb{C})$ 单连通。（见命题 21.34 与习题 7-16。）
\end{exercise}

\begin{exercise}\label{exercise:ch21-26}
由习题 21-25，$\operatorname{SL}(2,\mathbb{R})$ 微分同胚于 $\operatorname{SO}(2)\times\mathbb{R}^2\cong\mathbb{S}^1\times\mathbb{R}^2$，故其基本群同构于 $\mathbb{Z}$。令 $\widetilde{\operatorname{SL}}(2,\mathbb{R})$ 为 $\operatorname{SL}(2,\mathbb{R})$ 的万有覆叠群（定理 7.7）。证明 $\widetilde{\operatorname{SL}}(2,\mathbb{R})$ 不容许忠实表示，如下：设 $\rho\colon\widetilde{\operatorname{SL}}(2,\mathbb{R})\to\operatorname{GL}(V)$ 为任意表示。选 $V$ 的实基后可把 $\operatorname{GL}(V)$ 换成某 $\operatorname{GL}(n,\mathbb{R})$。则 $\rho$ 诱导 Lie 代数同态 $\rho_\ast\colon\mathfrak{sl}(2,\mathbb{R})\to\mathfrak{gl}(n,\mathbb{R})$。定义 $\varphi\colon\mathfrak{sl}(2,\mathbb{C})\to\mathfrak{gl}(n,\mathbb{C})$，$\varphi(A+\mathrm{i}B)=\rho_\ast A+\mathrm{i}\rho_\ast B$（$A,B\in\mathfrak{sl}(2,\mathbb{R})$）。
(a) 证明 $\varphi$ 是 Lie 代数同态。
(b) 证明存在 Lie 群同态 $\Phi\colon\operatorname{SL}(2,\mathbb{C})\to\operatorname{GL}(n,\mathbb{C})$ 使 $\Phi_\ast=\varphi$，且图表交换：$\widetilde{\operatorname{SL}}(2,\mathbb{R})\xrightarrow{\rho}\operatorname{GL}(n,\mathbb{R})$，经 $\operatorname{SL}(2,\mathbb{R})$ 到 $\operatorname{SL}(2,\mathbb{C})$ 的包含再经 $\Phi$ 到 $\operatorname{GL}(n,\mathbb{C})$。（提示：用习题 21-25。）
(c) 证明 $\rho$ 不忠实。
\end{exercise}
\end{document}
